%% ma: L12-B16-0011
%% lop: 12
%% chuong: 5
%% bai: 16
%% dang: 12.16.8
%% loai: TN4
%% muc: TH
%% dap_an: "B"
%% nguon: "Ảnh giáo viên gửi 04/11/2026 (Huỳnh Văn Ánh, Chuyên đề V - Phương pháp tọa độ trong không gian), Đề test theo bài, Phần I Câu 2"
%% kien_thuc: "Bài 16 (góc giữa đường thẳng và trục tọa độ), Bài 7 (vectơ đơn vị của trục), Bài 8 (tích vô hướng)"
%% trang_thai: cho-duyet
%% ly_do: "Dạng 12.16.8 mới đề xuất, chờ giáo viên duyệt."
%% ghi_chu: "Đáp án gốc B khớp (sympy). Lời giải gốc ghi nhầm vectơ chỉ phương là (1;2;2) trong khi đề cho (2;2;1); phép tính của lời giải gốc dùng (2;2;1) nên kết quả 2/3 vẫn đúng (với (1;2;2) sẽ ra 1/3)."
\begin{ex}
Trong không gian $Oxyz$, cho đường thẳng $d$ có vectơ chỉ phương $\vec u=(2;2;1)$. Gọi $\alpha$ là góc giữa $d$ và trục $Ox$. Tính $\cos\alpha$.
\choice{$\dfrac12$}{\True $\dfrac23$}{$\dfrac{2\sqrt2}{3}$}{$\dfrac13$}
\loigiai{Trục $Ox$ có vectơ chỉ phương $\vec i=(1;0;0)$. Khi đó
\[\cos\alpha=|\cos(\vec i,\vec u)|=\dfrac{|2\cdot1+2\cdot0+1\cdot0|}{\sqrt{1^2+0^2+0^2}\cdot\sqrt{2^2+2^2+1^2}}=\dfrac23.\]
Chọn B.}
\end{ex}
